\sect{ज्येष्ठ-कृष्णापरा-एकादशी-माहात्म्यम्}
\label{sec:padma-jyeshtha-krishnapara}

\ganapatyadiDhyanam
\padmaPuranam

\dnsub{अथ द्विपञ्चाशत्तमोऽध्यायः॥५२}

\uvacha{युधिष्ठिर उवाच}

\twolineshloka
{ज्येष्ठस्य कृष्णपक्षे तु किं नामैकादशी भवेत्}
{श्रोतुमिच्छामि माहात्म्यं तद् वदस्व जनार्दन}% ॥१॥

\uvacha{श्रीकृष्ण उवाच}

\twolineshloka
{साधु पृष्टं त्वया राजन् लोकानां हितकाम्यया}
{बहुपुण्यप्रदा ह्येषा महापातकनाशिनी}% ॥२॥

\twolineshloka
{अपरा नाम राजेन्द्र अपरा पुत्रदायिनी}
{लोके प्रसिद्धितां याति अपरां यस्तु सेवते}% ॥३॥

\twolineshloka
{ब्रह्महत्याभिभूतोऽपि गोत्रहा भ्रूणहा तथा}
{परापवादवादी च परस्त्रीरतिकोऽपि च}% ॥४॥

\twolineshloka
{अपरासेवनाद् राजन् विपाप्मा भवति ध्रुवम्}
{कूटसाक्ष्यं कूटमानं तुलाकूटं करोति यः}% ॥५॥

\twolineshloka
{कूटवेदं पठेद् यस्तु कूटशास्त्रं तथैव च}
{ज्योतिषी गणकः कूटः कूटायुर्वैदिको भिषक्}% ॥६॥

\twolineshloka
{कूटसाक्षिसमायुक्तो विज्ञेया नरकौकसः}
{अपरासेवनाद् राजन् पापैर्मुक्ता भवन्ति ते}% ॥७॥

\twolineshloka
{क्षत्रियः क्षात्रधर्मं यस्त्यक्त्वा युद्धात् पलायते}
{स याति नरकं घोरं स्वीयधर्मबहिष्कृतः}% ॥८॥

\twolineshloka
{अपरासेवनात् सोऽपि पापं त्यक्त्वा दिवं व्रजेत्}
{विद्यावान् यः स्वयं शिष्यो गुरुनिन्दां करोति च}% ॥९॥

\twolineshloka
{स महापातकैर्युक्तो निरयं याति दारुणम्}
{अपरासेवनात् सोऽपि सद्गतिं प्राप्नुयान्नरः}% ॥१०॥

\twolineshloka
{महिमानमपरायाः शृणु राजन् वदाम्यहम्}
{मकरस्थे रवौ माघे प्रयागे यत्फलं नृणाम्}% ॥११॥

\twolineshloka
{काश्यां यत् प्राप्यते पुण्यमुपरागे निमज्जनात्}
{गयायां पिण्डदानेन पितॄणां तृप्तिदो यथा}% ॥१२॥

\twolineshloka
{सिंहस्थिते देवगुरौ गौतम्यां स्नातको नरः}
{कन्यागते गुरौ राजन् कृष्णवेणीनिमज्जनात्}% ॥१३॥

\twolineshloka
{यत्फलं समवाप्नोति कुम्भकेदारदर्शनात्}
{बदर्याश्रमयात्रायां तत् तीर्थसेवनादपि}% ॥१४॥

\twolineshloka
{यत्फलं समवाप्नोति कुरुक्षेत्रे रविग्रहे}
{गजाश्वहेमदानेन यज्ञं कृत्वा सदक्षिणम्}% ॥१५॥

\twolineshloka
{तादृशं फलमाप्नोति अपराव्रतसेवनात्}
{अर्धप्रसूतां गां दत्त्वा सुवर्णं वसुधां तथा}% ॥१६॥

\twolineshloka
{नरो यत्फलमाप्नोति अपराया व्रतेन तत्}
{पापद्रुमकुठारीयं पापेन्धन दवानलः}% ॥१७॥

\twolineshloka
{पापान्धकारतरणिः पापसारङ्ग केसरी}
{बुद्बुदा इव तोयेषु पुत्तिका इव जन्तुषु}% ॥१८॥

\twolineshloka
{जायन्ते मरणायैव एकादश्या व्रतं विना}
{अपरां समुपोष्यैव पूजयित्वा त्रिविक्रमम्}% ॥१९॥

\threelineshloka
{सर्वपापविनिर्मुक्तो विष्णुलोके महीयते}
{लोकानां च हितार्थाय तवाग्रे कथितं मया}
{पठनाच्छ्रवणाद् राजन् गोसहस्रफलं लभेत्}% ॥२०॥

॥इति श्रीपाद्मे महापुराणे पञ्चपञ्चाशत्साहस्र्यां संहितायामुत्तरखण्डे उमापतिनारदसंवादान्तर्गतकृष्णयुधिष्ठिरसंवादे ज्येष्ठ-कृष्णापरा-एकादशी-माहात्म्यं नाम द्विपञ्चाशत्तमोऽध्यायः॥५२॥

\genMangalaShloka

\hyperref[sec:ekadashi_mahatmyam_padma_puranam]{\closesub}
\clearpage

